% Ad Familiares 1.6 --- headnote
% ad-familiares-01-06, March 56 BC, Rome

\textit{Cicero to P. Lentulus Spinther, proconsul of Cilicia,
written from Rome in March 56 BC. The shortest of the Lentulus
Spinther sequence (Fam.\ 1.1--1.6) and the closing pendant: the
Egyptian commission --- the question of whether Lentulus, as
governor of the province, would receive the senatorial charge
of restoring Ptolemy Auletes to the throne of Egypt --- has by
this date effectively collapsed. The 8 February riot at the
Milo contio (Fam.\ 1.5b, Q.\ fr.\ 2.3) has shown that Pompey
will not press the cause. The whole state of the affair Cicero
leaves to Pollio (Asinius Pollio, the future historian, then a
young man in his twenties, here in his earliest appearance in
the Ciceronian correspondence as a courier-cum-witness who has
``presided over'' the business at Rome).}

\textit{The body of the letter is consolation. Cicero's own
exile, returned from now four months, supplies the model: ``the
recollection of my own times consoles me, the image of which I
see in your affairs.'' The advice to ``show yourself the man
whom from your tender fingernails I have known'' borrows the
Greek tag --- \textit{ek hap\=al\=on onuch\=on} --- and the
closing assurance is in the formal style of patronage:
``every highest zeal and service from me; I shall not
disappoint your expectation.'' The letter is the bridge between
the dense early-56 BC sequence (Fam.\ 1.1--1.5b) and the rest of
the year, in which the Egyptian question recedes and Lentulus's
correspondent at Rome turns to other business: the Luca
compact, the post-Luca about-face, the great political-letter
arc of mid-56 BC.}
